\documentclass[10pt,a4paper]{amsart}
\usepackage[margin=19mm]{geometry}
\usepackage{amsmath,amssymb,mathtools}
\usepackage[scr=rsfs,cal=euler]{mathalfa}
\usepackage{xfrac}
\usepackage[unicode,hidelinks,bookmarks=false]{hyperref}
\newcommand{\ie}{i.e.\ }
\newcommand{\cf}{cf.\ }
\newcommand{\resp}{respectively\ }
\newcommand{\Subsection}{\subsection}
\providecommand{\nobreakdash}{\nobreak\hbox{-}}
\setlength{\emergencystretch}{2em}
\multlinegap0pt
\usepackage{mathtools}
\usepackage[normalem]{ulem}
\usepackage{amssymb}
\usepackage{enumerate}
\newtheorem{assumption}{Assumption}
\newtheorem{remark}{Remark}
\newtheorem{theorem}{Theorem}
\newtheorem{lemma}{Lemma}
\newcommand{\R}[0]{\mathbb R}
\newcommand{\V}[0]{\mathcal V}
\newcommand{\calH}[0]{\mathcal H}
\newcommand{\dd}[0]{\mathrm d}
\newcommand{\sumj}[0]{\sum_{j=1}^N}
\title{Erratum to: Port-Hamiltonian structure \\of interacting particle systems and its mean-field limit}
\author{Jannik Daun, Daniel Jannik Happ, Birgit Jacob, Claudia Totzeck}
\date{Source: arXiv:2301.06121v3 (CC BY 4.0)}
\begin{document}
\maketitle

\noindent\emph{Source and licence.} Erratum to: Port-Hamiltonian structure \\of interacting particle systems and its mean-field limit. Version arXiv:2301.06121v3 (CC BY 4.0), distributed by arXiv under the Creative Commons Attribution 4.0 International licence, http://creativecommons.org/licenses/by/4.0/, as recorded in the arXiv licence metadata for this exact version and shown on the abstract page https://arxiv.org/abs/2301.06121v3. Full licence notice: \texttt{licenses/arxiv-2301.06121-CC-BY-4.0.txt}.\par\smallskip

\noindent\textit{Editorial scope.} Counted unit: the article's lemma lem:invariantnew (source0.tex lines 1038--1043). The article's complete original proof of it is delivered with it (source0.tex lines 1044--1046, one \texttt{proof} environment of 3 lines). No mathematics is written, repaired or re-derived: the statement and the proof are the article's own lines, in the article's own notation, and no retained line is substituted or re-flowed.  Retained uncounted as the source-local context the proof needs: lines 836--844; lines 922--938; lines 970--972; lines 977--988.  Nothing was omitted from any retained span, so the package carries no omission record.  No retained line contains a citation, so the wrapper carries no bibliography and no bibliography entry.  No retained line contains a figure environment, an image-inclusion command or drawing code, so no image is delivered and the package declares no asset dependency.  Editorial formatting only, in addition: an AMS \texttt{amsart} preamble that restates the article's own declarations and macros used by the retained lines, namely the environments \texttt{assumption}, \texttt{remark}, \texttt{theorem}, \texttt{lemma} on separate counters, together with the standard packages those lines need.  The retained environments print in this document as Assumption 1, Remark 1, Theorem 1, Lemma 1 in order of appearance, whereas the article prints the same items with the numbering of its own full document.  The complete native source of the article as distributed in the arXiv e-print for this exact version is bundled unchanged as \texttt{sources/136106/source0.tex}.
\begin{assumption}\label{ass:particle}
\begin{itemize}
    \item [(1)] $F_i$ is continuous on $[0,T] \times \R^{dN} \times R^{dN}$ for all $i\in\{1,\dots,N\}$
    \item [(2)] For some $C>0$ it holds
    \[|F_i(t,x,v)| \le C(1+|x|+|v|) \quad \text{ for all }\quad i\in\{1,\dots,N\}, t\in[0,T] \text{ and } x,v\in \R^{dN}\]
    \item [(3)] For all $i\in\{1,\dots,N\}$, $F_i(t,x,v)$ is locally Lipschitz continuous w.r.t.~$x$ and $v$. In particular, for every compact set $K \subset \R^{dN} \times \R^{dN}$ there exists some $L_K>0$ such that 
    $F_i$ is on $K$ Lipschitz continuous  w.r.t.~$x$ and $v$ uniformly for all  $i\in\{1,\dots,N\}$.
\end{itemize}
\end{assumption}

\begin{remark}\label{rem:Psi}
For later use, we remark that  $\Psi(z)$ is diagonally dominant with non negative diagonal elements, hence positive semi-definite. Further, $\{ v \in \R^{Nd} : v = (\tilde v, \dots, \tilde v) \text{ for some } \tilde v \in \R^d \} $ is contained in the kernel of $\Psi(z)$. Moreover, if  $\psi(x)>0$  for $x\in\R$, then the eigenspace corresponding to the zero eigenvector is  spanned by $\{ v \in \R^{Nd} : v = (\tilde v, \dots, \tilde v) \text{ for some } \tilde v \in \R^d \} $.
\end{remark}
% 
As the positions of the particles appear only relatively in the dynamics, interacting particle systems are translational invariant w.r.t.~$x$. This motivates us to consider the dynamics with center of mass shifted to zero. Let us therefore  define the center of mass by $\bar x = \frac1N \sumj x_j$ and consider $z=(r,v)=(r_1,\dots,r_N,v_1,\dots,v_N)$ with $r_i = x_i - \bar x$ being the position relative to the center of mass. Note that the velocity of the center of mass is conserved. Indeed, using 
%$\psi(r) = \psi(-r)$ and 
$\nabla \mathcal V(-r) = -\nabla \mathcal V(r)$, we find
\begin{align}\label{eqn:meanv}
\sum_{i=1}^N\frac{\dd}{\dd t} v_i &= \frac{1}{N}\sum_{i=1}^N\sumj \psi(|r_j - r_i|)(v_j - v_i) - \frac{1}{N} \sum_{i=1}^N\sumj \nabla \V(r_i - r_j)=0,
\end{align}
Let us denote $\bar v := \frac1N \sumj v_j(0)$. The shifted dynamics is given by
\begin{subequations}\label{eq:phs}
\begin{align}\label{eqn:phs1}
    \frac{\dd}{\dd t} r_i &= v_i - \bar v,  &r_i(0) = (r_0)_i, \\
    \frac{\dd}{\dd t} v_i &= \frac{1}{N}\sumj \psi(| r_i - r_j |)(v_j - v_i) - \frac{1}{N} \sumj \nabla \mathcal V(r_ i - r_j),  &v_i(0)=(v_0)_i.\label{eqn:phs2}
\end{align}
\end{subequations}

\begin{equation}\label{eq:PHSparticle}
\frac{\dd}{\dd t} z = \left[ \begin{pmatrix} 0 & I \\ -I & 0 \end{pmatrix} - \begin{pmatrix} 0 & 0 \\ 0 & \Psi(z) \end{pmatrix} \right] \frac{\partial \calH^N(z)}{\partial z}. 
\end{equation}

\begin{theorem}\label{thm:existencenew}
Let Assumption~\ref{ass:particle}  hold, then for every initial condition $z_0=(r_0,v_0)\in \R^{Nd+Nd}$ the port-Hamiltonian system \eqref{eq:PHSparticle} possesses a unique global solution $z=(r,v)$ satisfying the dissipativity inequality
\begin{equation}\label{eqn:dissipnew}
\frac{\dd}{\dd t} \calH^N(z)= - \langle (v - \textbf{1}\bar v), \Psi(z) (v-\textbf{1}\bar v) \rangle \le 0,
\end{equation}
where 
\begin{equation*}
  \bar v:=\sum_{j=1}^N v_j(0). 
\end{equation*}
Further, the velocity of the center of mass is conserved, that is,
 $ \bar v=\sum_{j=1}^N v_j(t)$, $t\ge 0$. 
\end{theorem}

\begin{lemma}\label{lem:invariantnew}
Let Assumption~\ref{ass:particle} hold, $\psi(x)>0$  for $x\in\R$ and we define 
 \[ \mathcal N:= \{z=(r,v)\in X\mid  - \langle v - \mathbf{1}\bar v, \Psi(z)( v - \mathbf{1}\bar v )\rangle = 0\}. \]
 Let $z:[0,\infty)\rightarrow X$, $z=(r,v)$, be a solution of \eqref{eq:phs} with $z(t)\in \mathcal N$ for $t\ge 0$. Then the function $v$ is constant,
 \[  v_i(t)=\frac{1}{N}\sumj v_j(0)\quad\text{and}\quad \sumj \nabla \mathcal V(r_ i(t) - r_j(t))=0, \qquad i=1,\ldots, N, \, t\ge 0.\]
\end{lemma}

\begin{proof}
 As $\psi(x)>0$  for $x\in\R$,  we obtain $ v_i(t)=v_j(t)$ for  $i,j=1, \ldots, N$ and $t\ge 0$. This together with Theorem \ref{thm:existencenew} implies that $v$ is constant and satisfies $v_i(t)=\frac{1}{N}\sumj v_j(0)$ for $ i=1,\ldots, N$, and $t\ge 0$. The remaining statement follows directly from equation \eqref{eq:phs}. 
\end{proof}

\end{document}
